\documentclass[sigconf]{acmart}

% ACM KDD workshop — AI4Mental 2026
\acmConference[KDD AI4Mental '26]{Workshop on AI for Cognitive and Mental Health Support}{August 2026}{Jeju, South Korea}

\usepackage{booktabs}
\usepackage{multirow}
\usepackage{graphicx}
\usepackage{amsmath}
\usepackage{xcolor}
\usepackage{microtype}
\usepackage{enumitem}

\begin{document}

\title{Context-Dependent Facial Biomarkers of Depression:\\
From Clinical RCT to Cross-Corpus Generalization}

\author{[Authors]}
\affiliation{\institution{[Institution]}}
\email{[email]}

\begin{abstract}
Automated prediction of depression severity from facial dynamics is a long-standing goal
in affective computing, yet cross-corpus generalisation remains poorly understood.
We present a regression and classification study predicting depression indicators from
facial action units, gaze, and head pose (OpenFace~2.1.0) across two heterogeneous
corpora: the EmpkinS-EKSpression randomised controlled trial (RCT; $N=256$,
SCID-5-CV diagnosed, 4 conditions $\times$ 4 affective phases) and the Extended Distress
Analysis Interview Corpus (E-DAIC; $N=275$, unstructured clinical interview).
Within-corpus PHQ-8 regression achieves a 5-fold cross-validated CCC of 0.31 on
the best phase/condition (Preparatory/AFE), but cross-corpus transfer to E-DAIC collapses
to CCC\,$\approx 0$. In the reverse direction, modest signal (CCC\,$= 0.26$) emerges
only on the passive baseline phase---the RCT context most similar to an unstructured
interview. Cross-corpus binary classification shows a contrasting pattern: the
Preparatory/All configuration achieves AUC\,$= 0.70$ in forward transfer
(all Proposed participants as training), suggesting that categorical MDD status is more
transferable than continuous severity. We argue that \emph{functional context} is the
primary barrier to transfer and discuss implications for AI-based depression assessment.
\end{abstract}

\keywords{depression severity estimation, facial action units, cross-corpus generalisation,
affective computing, PHQ-8 regression, psychomotor dysregulation}

\maketitle

% =============================================================================
\section{Introduction}
% =============================================================================

Depression is one of the most prevalent and disabling mental health conditions worldwide,
affecting more than 280 million people and accounting for a substantial share of
years lived with disability~\cite{who2023}. Despite its prevalence, depression
remains substantially under-diagnosed: access to trained clinicians is limited,
stigma discourages help-seeking, and existing self-report instruments lack the
objectivity needed for population-scale monitoring. Automated analysis of
non-invasive behavioural signals---in particular, facial dynamics captured from
standard video---has emerged as a promising complementary approach for scalable,
objective depression screening~\cite{cummins2015review,pampouchidou2017automatic}.

Facial action units (AUs), defined by the Facial Action Coding System~\cite{ekman1978},
provide a fine-grained, theory-grounded representation of facial muscle movement.
In individuals with major depressive disorder (MDD), psychomotor dysregulation manifests
as blunted expressivity, reduced AU activity, prolonged lip-corner depression (AU15,
AU17), restricted gaze range, and decreased head-movement variability---patterns that
have been associated with depression severity in both clinical and community
samples~\cite{girard2014,jain2014depression,SchullerRingeval2015}. Competitive challenges
such as AVEC~\cite{ringeval2019avec} have pushed AU-based methods toward increasingly
high within-corpus performance, yet translating these results to unseen datasets remains
an open problem~\cite{cummins2015review,jaiswal2019mismatch}.

A fundamental and often underappreciated source of that generalisation gap is
\emph{elicitation context}. Psychomotor signs of depression are more pronounced
under structured affective challenge than under passive observation: studies using
standardised emotional tasks consistently report stronger facial biomarker effects
than those relying on naturalistic conversation or resting-state recordings, yet the
two paradigms are rarely compared within a single study and virtually never tested in
a cross-corpus transfer design~\cite{alghowinem2016multimodal,cummins2015review}.
This means that a model trained on structured-task data may fail not because of
domain shift in the feature distribution per se, but because the underlying
psychomotor signals are simply absent or attenuated in a naturalistic target corpus.

In our companion work~\cite{empkins_journal} on the EmpkinS-EKSpression randomised
controlled trial (RCT), we showed that facial AU features predict PHQ-8 depression
severity most reliably in the passive Preparatory baseline phase, not in the active
emotion regulation conditions---a counterintuitive result explained by the signal
being drowned by task-induced expression variability during active phases.
More broadly, that work characterised how the predictive value of facial features
varies across the four phases and four intervention conditions of the RCT.
The present paper asks the next logical question: do the most informative
phase/condition configurations from that RCT generalise to an entirely
independent corpus---the Extended Distress Analysis Interview Corpus
(E-DAIC~\cite{gratch2014edaic})---collected under an unstructured clinical
interview protocol with a different language and culture?

We run systematic bidirectional transfer experiments covering both continuous
PHQ-8 severity regression and binary MDD classification, and we show that
the degree of cross-corpus transfer depends critically on the functional alignment
between elicitation contexts rather than on the feature space or label instrument
alone. The passive Preparatory phase---the one RCT context most analogous to an
unstructured interview---is the only setting where meaningful regression transfer
occurs in the reverse direction, and Preparatory configurations also achieve the
strongest forward-transfer AUC (0.70). Importantly, binary classification
generalises more robustly than continuous regression, suggesting that categorical
MDD status is encoded in context-independent features (global expressivity level)
that continuous severity gradations are not.

\noindent\textbf{Contributions.}
\begin{itemize}[leftmargin=*,topsep=2pt,itemsep=1pt]
  \item The first systematic bidirectional cross-corpus study pairing a structured
        clinical RCT with a naturalistic interview corpus for both regression and
        classification.
  \item Empirical evidence that functional context alignment---not label mismatch or
        feature space---is the primary determinant of cross-corpus generalisation in
        facial depression assessment.
  \item A controlled comparison of label settings (SCID-5-CV vs.\ PHQ-8 threshold)
        showing that clinician-confirmed binary labels can compensate for nominal
        label-definition mismatch.
  \item Practical guidance: passive baseline observations transfer best; structured
        affective tasks should be incorporated into AI4Mental protocols targeting
        cross-context robustness.
\end{itemize}


% =============================================================================
\section{Related Work}
\label{sec:related}
% =============================================================================

\paragraph{Facial biomarkers of depression.}
Automated detection of depression from facial dynamics has been studied extensively
since the introduction of the AVEC benchmark challenges~\cite{ringeval2019avec,
valstar2016avec2016}. Early work demonstrated that AU-based features---particularly
AU15 (lip-corner puller), AU17 (chin raiser), and reduced AU12 (smile) activity---are
reliable markers of MDD~\cite{girard2014,jain2014depression,SchullerRingeval2015}.
Subsequent work extended these findings to continuous severity prediction,
showing that a combination of AU intensities, gaze direction, and head pose
can predict PHQ-8 or PHQ-9 scores with CCC values in the range of 0.3--0.6
within a single corpus~\cite{valstar2016avec2016,
ringeval2019avec,pampouchidou2017automatic}.
Deep learning approaches (3D-CNNs, attention-based temporal models) have pushed
within-corpus performance higher~\cite{yang2017dcnn,zhu2017multimodal},
and recent work has extended this to multimodal fusion with large language
models~\cite{anon2024multimodal,anon2023llm},
but the increased model capacity also increases sensitivity to corpus-specific
artefacts, making cross-corpus generalisation harder.

\paragraph{Cross-corpus and domain adaptation.}
Despite growing recognition of the generalisation problem, papers reporting
results on an unseen target corpus remain rare in the depression detection
literature. \citet{cummins2015review} surveyed the field and noted that the
overwhelming majority of studies report within-corpus leave-one-subject-out or
$k$-fold results, making it difficult to assess practical deployability.
\citet{jaiswal2019mismatch} directly examined cross-corpus AUC degradation for
audio-visual depression cues, finding drops of 15--30 percentage points when
transferring between interview corpora---even when the recording setup was similar.
Domain adaptation and adversarial feature alignment have been proposed to reduce
this gap~\cite{zhao2020emotion,zhang2020domain}, but these methods require
concurrent access to target-domain data, limiting applicability in real deployment.
Speech-based cross-corpus transfer between a German clinical interview corpus and
E-DAIC has recently been demonstrated~\cite{anon2025crosscultural}, showing that
MDD detection models can generalise across corpora and languages with acceptable
performance (F1\,=\,0.59 cross-corpus vs.\ 0.86 within-corpus).
To our knowledge, however, no prior work has examined cross-corpus transfer between
a structured experimental/RCT paradigm and a naturalistic interview corpus using
\emph{facial} features, which is the setting most relevant to clinical translation.

\paragraph{Contextual and cultural influences.}
Facial expression of affect is modulated by social display rules, cultural
background, and the specific elicitation context~\cite{matsumoto2008cultural,
alghowinem2016multimodal}. Several studies have noted that AU-based depression
markers are stronger and more consistent in structured paradigms
(standardised emotional stimuli, Trier Social Stress Test) than in
free-conversation settings~\cite{alghowinem2016multimodal,girard2014},
and facial features have also been shown to track affective states such as
stress reliably during structured emotion-evoking
tasks~\cite{rupp2025stress}.
Cross-language and cross-cultural transfer adds a further layer of difficulty:
German vs.\ English-speaking populations may exhibit different baseline
expressivity norms, affecting both the direction and magnitude of AU-based
depression indicators~\cite{matsumoto2008cultural}.
Our work contributes controlled evidence on all three axes: paradigm contrast
(structured RCT vs.\ interview), language (German vs.\ English), and
label protocol (SCID-5-CV vs.\ PHQ-8 self-report).

\paragraph{Label heterogeneity.}
Label mismatch is a frequently cited but rarely quantified source of cross-corpus
degradation. Different studies use PHQ-8, PHQ-9, BDI-II, HRSD, or clinician SCID
diagnoses as depression targets, and thresholds for binary classification vary
widely~\cite{cummins2015review}. \citet{ringeval2019avec} noted that PHQ-8 and
PHQ-9 scores are highly correlated but that converting continuous scores to
binary labels at different thresholds can substantially alter class balance and
apparent model performance. Our Setting~1 vs.\ Setting~2 comparison isolates
label mismatch from context mismatch, providing direct evidence on its
relative contribution.


% =============================================================================
\section{Datasets}
\label{sec:data}
% =============================================================================

\subsection{EmpkinS-EKSpression RCT (Proposed Corpus)}
\label{sec:empkins}

The EmpkinS-EKSpression study is a randomised controlled trial recruiting $N = 256$
participants ($n = 128$ MDD, $n = 128$ healthy controls), with MDD diagnosis confirmed
by the SCID-5-CV structured clinical interview~\cite{empkins_journal,anon2024protocol}.
Each participant was assigned to one of four intervention conditions: \textbf{AFE}
(Affect-Focused Expression), \textbf{CR} (Cognitive Restructuring), \textbf{CR+AFE}
(combined), and \textbf{SHAM} (active control).
Each session followed four sequential phases: a passive \textbf{Preparatory} baseline,
\textbf{Mood Induction} (standardised negative activation), \textbf{Negative Training}
(regulation with the assigned technique), and \textbf{Positive Training} (positive affect
elicitation and regulation). Video was recorded with a Sony camera and processed with
OpenFace~2.1.0~\cite{baltrusaitis2018openface}.
Depression severity was assessed with multiple scales; we use the
PHQ-8~\cite{kroenke2009phq8} as the regression target because it is the only scale
shared with E-DAIC. A methodologically important caveat is that the PHQ-8 was collected
\emph{remotely} (self-administered at home, before the lab visit), making it less
reliable than the in-lab, clinician-supervised scales (CES-D, HRSD variants) also
available in the dataset. The reliability ordering is SCID-5-CV\,$>$\,HRSD/CES-D\,$>$\,PHQ-8,
which must be kept in mind when interpreting regression performance.

\subsection{E-DAIC (Target Corpus)}
\label{sec:edaic}

The Extended Distress Analysis Interview Corpus (E-DAIC)~\cite{gratch2014edaic} comprises
$N = 275$ participants (66 depressed, 209 HC; 24\% positive rate), recorded during an
unstructured ${\sim}20$-minute semi-clinical interview conducted by a virtual agent.
Depression is defined by PHQ-8\,$\geq 10$. We use the official train+dev split ($n=219$)
for training and the held-out test set ($n=56$) for all reported evaluations.
Features were extracted with OpenFace~2.1.0, matching the Proposed corpus pipeline.
The two corpora share the same feature extractor and label instrument (PHQ-8), but differ
on every other axis: diagnosis protocol (SCID-5-CV vs.\ self-report threshold),
elicitation context (structured 4-phase RCT vs.\ unstructured interview),
language and culture (German vs.\ English-speaking US), and class balance (50/50 vs.\ 24\%).
These contrasts are the primary subject of the discussion in Section~\ref{sec:discussion}.


% =============================================================================
\section{Methods}
\label{sec:methods}
% =============================================================================

\subsection{Feature Extraction and Alignment}
\label{sec:features}

Both corpora were processed with OpenFace~2.1.0~\cite{baltrusaitis2018openface}, yielding
per-frame estimates of facial action unit intensities and presence indicators (18 AUs),
bilateral gaze direction vectors, gaze angles, and head pose (translation and rotation).
Low-quality frames (OpenFace confidence\,$< 0.5$) were discarded; participants with fewer
than 50 retained frames were excluded.
Each signal column was summarised into a fixed-length feature vector using 18 statistical
functionals applied to both the original signal and its first-order frame differences:
mean, std, min, max, skew, kurtosis, range, and Shannon entropy for each; plus
rate-of-change and peak count on the original signal.
After applying this aggregation to both corpora, we intersected their column sets to obtain
882 shared features covering all AUs, gaze angles, and head pose rotations, while excluding
corpus-specific landmark positions and composite expressions.
All cross-corpus experiments use this aligned 882-feature set.

\subsection{Machine Learning Pipeline}
\label{sec:pipeline}

All preprocessing steps are fit exclusively on training data and applied without
refitting to the test set, ensuring no data leakage.
For the Proposed corpus we use pre-computed stratified 80/20 splits (stratified by
diagnosis and condition), shared across all phase/condition combinations to ensure
comparability in cross-corpus settings.
For E-DAIC we use the official train+dev\,/\,test partition.
Missing values, arising from participants excluded in a given phase, are imputed with the
training-set column median.

Feature selection proceeds in two steps, both computed on the training set only.
First, features are ranked by absolute Spearman correlation with the target; those
surviving Benjamini--Hochberg FDR correction ($q = 0.05$) are retained (with a fallback
of the top 20 by $|\rho|$ when fewer than 20 pass the threshold).
Second, recursive feature elimination (RFE) with a Ridge surrogate model further removes
redundant features.
For regression we evaluate 12 models: Ridge, Lasso, ElasticNet, SVR (RBF), $k$-NN,
Random Forest, Extra Trees, Gradient Boosting, XGBoost~\cite{chen2016xgboost},
LightGBM~\cite{ke2017lightgbm}, CatBoost~\cite{prokhorenkova2018catboost}, and MLP.
For classification we replace these with Logistic Regression, SVC, and the same
ensemble and neural models.
Three scalers (Standard, MinMax, Robust) are evaluated jointly with each model.
Hyperparameters are selected via 5-fold grid search on the training set.

We report CCC (Lin's concordance correlation coefficient~\cite{lin1989ccc}) as the primary
regression metric, as it jointly penalises poor rank-order correlation and systematic
scale or offset bias---a property particularly relevant for cross-corpus evaluation
where the latter is common. Pearson $r$, MAE, and RMSE are reported as secondary metrics.
For classification, AUC is the primary metric (threshold-independent); macro-averaged
F1 (F1$_{\rm M}$\,=\,mean of per-class F1) is reported as secondary, as it treats both
classes symmetrically and is better suited to imbalanced test sets than binary F1 for
the depressed class alone.

\subsection{Experimental Designs}
\label{sec:experiments}

\subsubsection{Within-Corpus Regression (Proposed Corpus)}
\label{sec:within}

To comprehensively characterise which phase/condition combinations carry the most
PHQ-8 signal, we apply 5-fold stratified nested cross-validation (outer CV for evaluation,
inner CV for hyperparameter search) across all 4 phases $\times$ 5 conditions (AFE, CR,
CR+AFE, SHAM, and all conditions pooled), yielding 20 combinations in total.
CCC reported in Table~\ref{tab:within} are outer-CV means across folds.
We additionally evaluate four clinically relevant scales (PHQ-8, PHQ-9, CES-D, HRSD-17)
to characterise target reliability.

\subsubsection{Cross-Corpus Regression (Proposed $\leftrightarrow$ E-DAIC)}
\label{sec:cross}

We design two directional transfer experiments using the same aligned 882-feature set.
\textbf{Experiment~C1} trains on the Proposed 80\% training split and tests on the
E-DAIC official test set ($n=56$), assessing forward transfer from the RCT to the
clinical interview.
\textbf{Experiment~B} reverses the direction: it trains on E-DAIC train+dev ($n=219$)
and tests on the Proposed 20\% holdout, asking whether interview-trained models
generalise to RCT phases.
In both experiments, preprocessing and feature selection are fit on the source partition
only and applied without refitting to the target.
Within-corpus baselines (\textbf{D1}: Proposed\,$\to$\,Proposed; \textbf{D2}: E-DAIC\,$\to$\,E-DAIC)
are trained and tested under the same fixed splits to anchor performance.

Experiments use the seven source configurations in Table~\ref{tab:configs}, which cover
the three latency-phase variants and four phase/condition combinations identified as
most discriminative in within-corpus analysis.

\subsubsection{Cross-Corpus Classification (Proposed $\leftrightarrow$ E-DAIC)}
\label{sec:cross_cls}

The same bidirectional design (B, C1) and within-corpus baselines (D1, D2) are applied
to binary depression classification (MDD vs.\ HC).
We additionally include \textbf{Experiment~A}: training on \emph{all} Proposed
participants (no holdout, $n \approx 60$--$251$ depending on condition) and testing on
the E-DAIC official test set---this exploits the maximum available RCT training data
when no within-corpus test set is required.

We evaluate two label settings to isolate the effect of label mismatch.
\textbf{Setting~1} uses SCID-5-CV binary labels for the Proposed corpus and PHQ-8\,$\geq 10$
for E-DAIC (diagnostic protocol mismatch).
\textbf{Setting~2} uses PHQ-8\,$\geq 10$ for both corpora (matched label definition).

The same seven source configurations (Table~\ref{tab:configs}) are used for both
regression and classification experiments.

\begin{table}[t]
\small
\caption{Source configurations for cross-corpus experiments, ranked by within-corpus
PHQ-8 CCC (5-fold nested CV mean\,$\pm$\,SD). All three latency configurations and
the four best phase/condition combinations are included.}
\label{tab:configs}
\renewcommand{\arraystretch}{1.05}
\begin{tabular}{llll}
\toprule
Short name & Phase & Condition & PHQ-8 CCC \\
\midrule
Prep./AFE      & Preparatory       & AFE          & $0.31 \pm 0.22$ \\
Prep./SHAM     & Preparatory       & SHAM         & $0.28 \pm 0.17$ \\
Prep./All      & Preparatory       & All (pooled) & $0.04 \pm 0.15$ \\
Pos./AFE      & Positive Training & AFE          & $0.26 \pm 0.33$ \\
Neg./CR+AFE   & Negative Training & CR+AFE       & $0.09 \pm 0.26$ \\
Neg./CR       & Negative Training & CR           & $0.20 \pm 0.20$ \\
Neg./SHAM     & Negative Training & SHAM         & $0.26 \pm 0.16$ \\
\bottomrule
\multicolumn{4}{l}{\footnotesize Prep./All includes all 4 conditions pooled; largest training set.}
\end{tabular}
\end{table}


% =============================================================================
\section{Results}
\label{sec:results}
% =============================================================================

\subsection{Within-Corpus Regression (Proposed Corpus)}
\label{sec:results_within}

Table~\ref{tab:within} summarises 5-fold nested CV performance across all 20
phase/condition combinations and four depression scales. PHQ-8 CCC peaks at 0.31
(Preparatory/AFE) and 0.28 (Preparatory/SHAM), confirming that the passive baseline
phase carries the strongest and most reliable PHQ-8 signal in this RCT. Active
training phases (Negative, Positive) achieve CCCs up to 0.26--0.28 in specific
conditions, while the pooled Preparatory/All setting drops to CCC\,=\,0.04, suggesting
that condition-specific dynamics are more informative than size.
PHQ-9 and CES-D generally track PHQ-8, while HRSD-17 shows weaker and noisier CCCs,
consistent with its rating-scale format being harder to predict from passive facial dynamics.
The low absolute CCC values across all cells (maximum 0.36, CES-D/Pos./SHAM) reflect
both the limited reliability of the remotely collected PHQ-8 and the conservative
nature of nested CV compared to single holdout estimates.

\begin{table*}[t]
\centering
\footnotesize
\renewcommand{\arraystretch}{1.05}
\caption{Within-corpus PHQ-8 regression (Proposed corpus, 882-feature restricted set).
5-fold stratified nested CV. CCC = mean across outer folds; (SD) in parentheses for PHQ-8.
MAE reported for PHQ-8 only. Bold = best PHQ-8 CCC within each phase group.}
\label{tab:within}
\resizebox{\textwidth}{!}{%
\begin{tabular}{ll rr rrr}
\toprule
& & \multicolumn{2}{c}{\textbf{PHQ-8}} & \textbf{PHQ-9} & \textbf{CES-D} & \textbf{HRSD-17} \\
\cmidrule(lr){3-4}
Phase & Condition & CCC\,(SD) & MAE & CCC & CCC & CCC \\
\midrule
\multirow{5}{*}{Preparatory}
 & AFE       & \textbf{0.31}\,(0.22) & 5.34 & 0.34 & 0.17 & 0.06 \\
 & CR+AFE    & 0.17\,(0.11) & 5.03 & 0.03 & 0.03 & --0.10 \\
 & CR        & 0.00\,(0.20) & 5.83 & 0.34 & 0.07 & 0.22 \\
 & SHAM      & 0.28\,(0.17) & 4.85 & 0.22 & 0.35 & 0.18 \\
 & All cond. & 0.04\,(0.15) & 5.11 & 0.13 & 0.14 & 0.10 \\
\midrule
\multirow{5}{*}{Mood Induction}
 & AFE       & --0.01\,(0.08) & 7.34 & --0.01 & 0.09 & 0.01 \\
 & CR+AFE    & 0.02\,(0.15) & 5.54 & 0.13 & 0.10 & 0.24 \\
 & CR        & 0.09\,(0.24) & 5.18 & 0.00 & 0.16 & 0.12 \\
 & SHAM      & 0.03\,(0.12) & 6.05 & 0.17 & 0.03 & 0.10 \\
 & All cond. & \textbf{0.14}\,(0.05) & 5.10 & 0.12 & 0.11 & 0.05 \\
\midrule
\multirow{5}{*}{Negative Training}
 & AFE       & 0.19\,(0.18) & 5.14 & 0.20 & 0.13 & --0.03 \\
 & CR+AFE    & 0.09\,(0.26) & 4.86 & 0.07 & --0.03 & 0.18 \\
 & CR        & 0.20\,(0.20) & 5.11 & 0.20 & 0.18 & 0.08 \\
 & SHAM      & \textbf{0.26}\,(0.16) & 5.03 & 0.28 & 0.20 & 0.25 \\
 & All cond. & 0.13\,(0.10) & 5.26 & 0.11 & 0.07 & 0.07 \\
\midrule
\multirow{5}{*}{Positive Training}
 & AFE       & 0.26\,(0.33) & 5.19 & --0.11 & --0.08 & 0.12 \\
 & CR+AFE    & --0.04\,(0.14) & 5.93 & 0.17 & 0.07 & 0.29 \\
 & CR        & 0.18\,(0.24) & 5.90 & 0.21 & 0.20 & 0.19 \\
 & SHAM      & \textbf{0.24}\,(0.23) & 4.87 & 0.33 & 0.36 & 0.26 \\
 & All cond. & 0.09\,(0.05) & 5.08 & 0.14 & 0.05 & 0.10 \\
\bottomrule
\end{tabular}}
\end{table*}

\subsection{Cross-Corpus Regression (Proposed $\leftrightarrow$ E-DAIC)}
\label{sec:results_cross}

Table~\ref{tab:cc_reg} reports CCC and MAE for all seven configurations.
Using all Proposed participants for training (Exp.~A, Proposed$_{\rm all}$\,$\to$\,E-DAIC)
does not improve forward transfer: CCC reaches at most 0.15 (Neg./CR), confirming that
the near-zero forward performance is driven by contextual mismatch, not training set size.
Restricting to the training split only (Exp.~C1) gives similar results: forward transfer
collapses almost entirely, with CCC\,$\approx 0$ for all latency configurations and
a modest exception for Neg./CR (CCC\,=\,0.21), possibly reflecting a chance alignment
between a specific negative-training feature pattern and interview-context dysregulation.
The reverse direction (E-DAIC\,$\to$\,Proposed, Exp.~B) yields more consistent signal:
CCC\,=\,0.26 for Prep./SHAM and CCC\,=\,0.25 for Neg./CR+AFE, while the pooled
Prep./All drops to 0.10. The within-corpus Proposed baseline (D1) peaks at CCC\,=\,0.28
(Neg./CR), confirming that even within-corpus PHQ-8 regression is difficult under
the restricted 882-feature set.
The E-DAIC within-corpus baseline is config-independent (fixed split) at CCC\,=\,0.08.

\begin{table*}[t]
\centering
\footnotesize
\renewcommand{\arraystretch}{1.1}
\caption{Cross-corpus PHQ-8 regression. Bold = best CCC within each experiment row.
$\dagger$~E-DAIC within-corpus result is config-independent (fixed train+dev/test split).}
\label{tab:cc_reg}
\resizebox{\textwidth}{!}{%
\begin{tabular}{l rr rr rr rr rr rr rr}
\toprule
& \multicolumn{6}{c}{Preparatory phase} & \multicolumn{8}{c}{Best training phase per condition} \\
\cmidrule(lr){2-7}\cmidrule(lr){8-15}
& \multicolumn{2}{c}{\textbf{Prep./AFE}} & \multicolumn{2}{c}{\textbf{Prep./SHAM}}
& \multicolumn{2}{c}{\textbf{Prep./All}} & \multicolumn{2}{c}{\textbf{Pos./AFE}}
& \multicolumn{2}{c}{\textbf{Neg./CR+AFE}} & \multicolumn{2}{c}{\textbf{Neg./CR}}
& \multicolumn{2}{c}{\textbf{Neg./SHAM}} \\
\cmidrule(lr){2-3}\cmidrule(lr){4-5}\cmidrule(lr){6-7}
\cmidrule(lr){8-9}\cmidrule(lr){10-11}\cmidrule(lr){12-13}\cmidrule(lr){14-15}
& CCC & MAE & CCC & MAE & CCC & MAE & CCC & MAE & CCC & MAE & CCC & MAE & CCC & MAE \\
\midrule
\multicolumn{15}{c}{\textbf{Cross-corpus}} \\
\midrule
Proposed$_{\rm all}$ $\to$ E-DAIC
  & 0.03 & 6.13 & 0.00 & 5.83 & 0.00 & 7.18
  & --0.03 & 6.86 & 0.13 & 5.62 & \textbf{0.15} & 5.77 & 0.10 & 5.49 \\
Proposed $\to$ E-DAIC
  & 0.01 & 6.98 & 0.08 & 5.40 & 0.03 & 9.59
  & 0.04 & 8.76 & 0.03 & 8.75 & \textbf{0.21} & 8.87 & 0.03 & 6.27 \\
E-DAIC $\to$ Proposed
  & 0.14 & 6.83 & \textbf{0.26} & 5.50 & 0.10 & 4.77
  & 0.10 & 7.25 & 0.25 & 5.40 & 0.04 & 5.31 & 0.14 & 5.06 \\
\midrule
\multicolumn{15}{c}{\textbf{Within-corpus baseline}} \\
\midrule
Proposed $\to$ Proposed
  & 0.03 & 6.42 & 0.19 & 6.40 & 0.09 & 4.63
  & 0.10 & 6.00 & 0.05 & 6.25 & \textbf{0.28} & 3.91 & 0.14 & 4.78 \\
E-DAIC $\to$ E-DAIC$\,\dagger$
  & 0.08 & 5.39 & 0.08 & 5.39 & 0.08 & 5.39
  & 0.08 & 5.39 & 0.08 & 5.39 & 0.08 & 5.39 & 0.08 & 5.39 \\
\bottomrule
\end{tabular}}
\end{table*}

\subsection{Cross-Corpus Classification (Proposed $\leftrightarrow$ E-DAIC)}
\label{sec:results_cls}

Table~\ref{tab:cc_cls} reports AUC and macro F1 (F1$_{\rm M}$) for both label settings.
In Setting~1 (SCID vs.\ PHQ-8 label mismatch), Experiment~A with Prep./All achieves
AUC\,=\,0.70---the strongest forward-transfer result across both tasks---likely because
pooling all four conditions maximises the Proposed training set size ($n=251$) while
the passive Preparatory phase aligns best with the E-DAIC interview context.
With the Proposed train split only (Exp.~C1), forward AUC peaks at 0.65 (Neg./CR).
Reverse transfer (E-DAIC\,$\to$\,Proposed, Exp.~B) achieves high AUC on small individual
condition test sets, including 1.00 for Pos./AFE ($n_\text{test}=12$; interpret with caution).
In Setting~2 (matched PHQ-8 labels), forward AUC is marginally lower (best: 0.67,
Neg./SHAM in Exp.~A), while reverse transfer peaks at 0.85 (Neg./SHAM, $n_\text{test}=13$).
The within-corpus Proposed baseline (D1) reaches AUC\,=\,0.78 (Neg./CR, Setting~1)
and 0.91 (Neg./CR, Setting~2; $n_\text{test}=12$), while the E-DAIC within-corpus
baseline is AUC\,=\,0.58 (config-independent).

\begin{table*}[t]
\centering
\footnotesize
\renewcommand{\arraystretch}{1.1}
\caption{Cross-corpus binary depression classification. AUC (primary) and macro-averaged
F1 (F1$_{\rm M}$; mean of per-class F1) per configuration and label setting.
Bold = best AUC within each experiment row.
$\dagger$~E-DAIC within-corpus result is config-independent (fixed train+dev/test split).
$^*$~$n_\mathrm{test} \le 13$; interpret with caution.}
\label{tab:cc_cls}
\resizebox{\textwidth}{!}{%
\begin{tabular}{l rr rr rr rr rr rr rr}
\toprule
& \multicolumn{6}{c}{Preparatory phase} & \multicolumn{8}{c}{Best training phase per condition} \\
\cmidrule(lr){2-7}\cmidrule(lr){8-15}
& \multicolumn{2}{c}{\textbf{Prep./AFE}} & \multicolumn{2}{c}{\textbf{Prep./SHAM}}
& \multicolumn{2}{c}{\textbf{Prep./All}} & \multicolumn{2}{c}{\textbf{Pos./AFE}}
& \multicolumn{2}{c}{\textbf{Neg./CR+AFE}} & \multicolumn{2}{c}{\textbf{Neg./CR}}
& \multicolumn{2}{c}{\textbf{Neg./SHAM}} \\
\cmidrule(lr){2-3}\cmidrule(lr){4-5}\cmidrule(lr){6-7}
\cmidrule(lr){8-9}\cmidrule(lr){10-11}\cmidrule(lr){12-13}\cmidrule(lr){14-15}
& AUC & F1$_{\rm M}$ & AUC & F1$_{\rm M}$ & AUC & F1$_{\rm M}$ & AUC & F1$_{\rm M}$ & AUC & F1$_{\rm M}$ & AUC & F1$_{\rm M}$ & AUC & F1$_{\rm M}$ \\
\midrule
\multicolumn{15}{c}{\textbf{Setting 1: Proposed SCID labels --- E-DAIC PHQ-8$\,\geq 10$}} \\
\midrule
Proposed$_{\rm all}$ $\to$ E-DAIC
  & 0.62 & 0.30 & 0.60 & 0.50 & \textbf{0.70} & 0.59
  & 0.55 & 0.28 & 0.67 & 0.42 & 0.57 & 0.31 & 0.66 & 0.41 \\
Proposed $\to$ E-DAIC
  & 0.59 & 0.50 & 0.55 & 0.41 & 0.57 & 0.23
  & 0.58 & 0.58 & 0.57 & 0.48 & \textbf{0.65} & 0.46 & 0.59 & 0.41 \\
E-DAIC $\to$ Proposed
  & 0.88 & 0.33 & 0.79 & 0.58 & 0.57 & 0.52
  & \textbf{1.00}$^*$ & 0.33 & 0.67 & 0.50 & 0.69 & 0.44 & 0.74 & 0.58 \\
Proposed $\to$ Proposed
  & 0.67 & 0.33 & 0.36 & 0.31 & 0.66 & 0.67
  & 0.47 & 0.50 & 0.47 & 0.56 & \textbf{0.78} & 0.62 & 0.64 & 0.45 \\
E-DAIC $\to$ E-DAIC$\,\dagger$
  & 0.58 & 0.50 & 0.58 & 0.50 & 0.58 & 0.50
  & 0.58 & 0.50 & 0.58 & 0.50 & 0.58 & 0.50 & 0.58 & 0.50 \\
\midrule
\multicolumn{15}{c}{\textbf{Setting 2: PHQ-8$\,\geq 10$ labels for both corpora}} \\
\midrule
Proposed$_{\rm all}$ $\to$ E-DAIC
  & 0.57 & 0.23 & \textbf{0.69} & 0.35 & 0.66 & 0.23
  & 0.62 & 0.23 & 0.61 & 0.59 & 0.64 & 0.29 & 0.67 & 0.29 \\
Proposed $\to$ E-DAIC
  & 0.61 & 0.39 & 0.53 & 0.50 & \textbf{0.63} & 0.30
  & 0.59 & 0.37 & 0.58 & 0.45 & 0.57 & 0.23 & 0.54 & 0.46 \\
E-DAIC $\to$ Proposed
  & 0.78 & 0.50 & 0.67 & 0.19 & 0.57 & 0.30
  & 0.78 & 0.25 & 0.77 & 0.69 & 0.63 & 0.49 & \textbf{0.85}$^*$ & 0.68 \\
Proposed $\to$ Proposed
  & 0.48 & 0.50 & 0.83 & 0.46 & 0.60 & 0.65
  & 0.63 & 0.56 & 0.60 & 0.31 & \textbf{0.91}$^*$ & 0.70 & 0.57 & 0.41 \\
E-DAIC $\to$ E-DAIC$\,\dagger$
  & 0.58 & 0.50 & 0.58 & 0.50 & 0.58 & 0.50
  & 0.58 & 0.50 & 0.58 & 0.50 & 0.58 & 0.50 & 0.58 & 0.50 \\
\bottomrule
\end{tabular}}
\end{table*}

% =============================================================================
\section{Discussion}
\label{sec:discussion}
% =============================================================================

The most striking finding is the near-complete collapse of regression performance when
transferring Proposed-trained models to E-DAIC (CCC\,$\approx 0$, Exp.~C1),
despite both corpora sharing the same feature extractor and label instrument.
Since both PHQ-8 measurements are self-administered, label format alone cannot explain
the failure. The primary cause is instead \emph{contextual mismatch}: the facial dynamics
elicited during structured affective tasks---erratic AU14 intensity peaks, restricted
gaze entropy, reduced head pitch variability---are signatures of psychomotor dysregulation
under challenge, and these dynamics are not present during an unstructured interview.
The modest reverse-direction signal (E-DAIC\,$\to$\,Proposed, CCC\,=\,0.26 for Prep./SHAM,
Pearson $r \approx 0.45$) is consistent with this interpretation: the passive
Proposed \emph{Preparatory} phase is functionally the closest to an unstructured interview,
and it is the only phase where interview-trained models find transferable structure.

The cross-corpus classification results tell a more nuanced story.
The Prep./All configuration (all participants, Preparatory phase) achieves AUC\,=\,0.70
in Experiment~A (all Proposed\,$\to$\,E-DAIC test, Setting~1)---the strongest forward-transfer
result across both tasks. This result has two drivers: first, the Preparatory phase is the
Proposed context most functionally aligned with an unstructured interview; second,
pooling all four conditions maximises the Proposed training set ($n=251$), providing
enough data for a stable classifier despite the domain shift.
Notably, matched labels (Setting~2, PHQ-8 for both) do not consistently outperform
Setting~1 (SCID labels), suggesting that the high reliability of the SCID-5-CV
clinician diagnosis may compensate for the nominal label mismatch.
We attribute the greater transferability of binary classification compared to
continuous regression to the nature of the respective targets: categorical MDD status
may be encoded in more context-robust facial features (e.g.\ overall AU activity level,
global expressivity reduction), whereas PHQ-8 severity estimation requires fine-grained
gradations that are highly sensitive to the elicitation context.

Comparing within-corpus to cross-corpus performance reveals a consistent domain gap.
In the Proposed$\,\to\,$E-DAIC direction, all seven configurations show a decline from
the within-corpus Proposed baseline (D1), with Neg./CR showing the smallest drop
(0.78\,$\to$\,0.65) and Prep./SHAM already weak within-corpus (0.36) yet improving
cross-corpus (0.55).
In the reverse direction, starting from the config-independent E-DAIC baseline
(AUC\,=\,0.58), cross-corpus AUC varies widely (0.57--1.00), with the highest values
concentrated on small test sets ($n \le 13$).
Across both directions, Preparatory-phase configurations achieve the most consistent
cross-corpus generalisation, reinforcing the role of functional context alignment.

Extending training and test sets to the full corpora without held-out splits (all
Proposed participants\,$\to$\,all E-DAIC, and vice versa) did not improve results
over split-based evaluations (regression CCC\,$\leq 0.21$, classification
AUC\,$\leq 0.70$), confirming that domain mismatch rather than limited training size
is the primary barrier to cross-corpus generalisation.

These findings carry a direct implication for AI4Mental system design: passive observation
and unstructured interviews are unlikely to elicit the facial biomarkers necessary for
robust severity estimation. Systems targeting generalisation should incorporate structured
affective tasks as part of the assessment protocol, and evaluation should distinguish
between coarse binary detection and continuous severity prediction.


% =============================================================================
\section{Conclusion}
% =============================================================================

We presented a cross-corpus study of facial depression assessment pairing a structured
RCT (EmpkinS-EKSpression, Proposed corpus) with a naturalistic clinical interview corpus
(E-DAIC). PHQ-8 regression transfer fails almost entirely in the forward direction
(CCC\,$\approx 0$), while binary MDD classification transfers more robustly
(AUC\,=\,0.70, Prep./All, Setting~1). Reverse regression transfer succeeds modestly
only on the passive Preparatory baseline phase, the one RCT context functionally comparable
to an interview. Together, the results establish functional context as the primary
determinant of cross-corpus generalisation in facial depression assessment and argue
for structured affective task designs in AI-based mental health tools.


% =============================================================================
\bibliographystyle{ACM-Reference-Format}
\bibliography{references}

\end{document}
